\documentclass[11pt]{article}
\usepackage{mathunal-base}
\mathunalfuente{serif}

\renewcommand{\examtitulo}{Primer Parcial de Matemáticas III}
\renewcommand{\examfecha}{Semestre 02 de 2007 --- Septiembre 24 de 2007}

\begin{document}

\examheader

\vspace{4pt}
\noindent Nombre: \dotfill\ Profesor: \dotfill

\begin{center}\textbf{Duración: 1 hora y 50 minutos}\end{center}

\noindent\textbf{Observaciones:} - Justifique cuidadosamente sus respuestas.

\hspace{2.3cm} - Cualquier tentativa de fraude es causal de anulación del examen.

\vspace{8pt}
\begin{enumerate}
\item{} (25\%) Sea
\[
f(x,y) = \left\{
\begin{array}{ll}
\dfrac{y^4}{x^2+y^2}+x, & (x,y)\neq(0,0)\\[6pt]
0, & (x,y)=(0,0).
\end{array}
\right.
\]
\noindent Muestre que $f$ es diferenciable en $(0,0)$.

\vspace{5cm}

\item{} (25\%)
\begin{enumerate}
\item{} (10\%) Halle la curva de nivel de valor 1 de la función $f(x,y)=\dfrac{y+1}{x^2+1}$ y esbócela en el plano $\mathbb{R}^2$.

\vspace{4cm}

\item{} (15\%) Determine si existe o no
\[
\lim_{(x,y)\to(0,0)}\left(\frac{xy}{x^4+y^2}-3\right).
\]
\noindent (Sugerencia: Ensaye primero algunas trayectorias adecuadas).

\vspace{4cm}
\end{enumerate}

\item{} (25\%) Calcule el valor de la mayor derivada direccional en $(1,1)$ de la función $f$ definida por $f(x,y)=g(x^2+2y^2,xy)$ si se sabe que
\begin{enumerate}
\item{} $\nabla g(1,1) = (3,1)$
\item{} $\nabla g(3,1) = (0,3)$.
\end{enumerate}
\noindent (Observación: Para resolver este problema es posible que no se requiera usar toda la información dada sobre $g$).

\vspace{6cm}

\item{} (25\%) Determine el punto que está en el paraboloide $x^2+y^2+z=9$ y en el cual la recta normal a esta superficie es perpendicular a los vectores normales de los planos $y+z=-16$ y $x-y+z=4$.

\noindent (sugerencia: Establezca la relación entre un vector director de la recta normal y los vectores normales a los planos).

\vspace{6cm}
\end{enumerate}

\examfooter
\end{document}
